% Shared body of the Phase F2 book fixture (report.tex, report-editorial.tex).
% Fictional content: the Tessaly Valley Energy Cooperative, its sites and
% every figure below are invented for this demonstration. The body uses only
% book, long-form and shared semantic primitives -- no font, color or
% spacing commands -- so the selected theme alone decides how it looks.

\RKFrontMatterBegin
\RKTitlePage[A field handbook for community microgrids]{The Quiet Grid}

\begin{bookdetails}{The Quiet Grid: a field handbook for community microgrids}
\bookdetail{Edition}{First edition, March 2027. Prepared by the fictional Tessaly
Valley Energy Cooperative as a demonstration of the ReportKit book publication type.}
\bookdetail{Licence}{Original prose in this demonstration is licensed under
\href{https://creativecommons.org/licenses/by/4.0/}{CC BY 4.0}. The cooperative, its
sites, people and measurements are fictional.}
\bookdetail{Disclaimer}{This handbook is not engineering advice. Every figure in it is
invented to exercise the book structure: part and chapter openers, appendices, a
glossary and a reference list.}
\bookdetail{Source}{\RKPath{https://example.org/quiet-grid-handbook}}
\end{bookdetails}

\RKContents
\RKMainMatterBegin

\bookpart{Foundations}

\section{Why small grids fail quietly}
\label{ch:quiet-failure}

A community microgrid rarely fails with a bang. It fails with a slow drift: a battery
bank that holds a little less each winter, an inverter that trips once a month and is
reset by whoever notices first, a load that nobody measured because it arrived after
commissioning. The cooperative's first three sites all ran for more than a year before
anyone called them unreliable, and all three had been unreliable for most of that year.

This handbook is written for the people who inherit those systems: the volunteer
operator, the district engineer who visits twice a year, and the committee that must
decide whether to extend a grid or replace it. It assumes no specialist training, only
the patience to measure before deciding \cite{valley-survey}.

\subsection{Three kinds of quiet failure}

The failures we saw fall into three groups. \emph{Capacity drift} is the gradual gap
between what a system was sized for and what it now serves. \emph{Maintenance debt} is
the accumulation of small deferred tasks, each harmless alone. \emph{Knowledge loss}
happens when the one person who understood the system moves away and nobody wrote down
what they knew.

\begin{principle}{Measure the load you have, not the load you planned}
Every sizing decision in this handbook starts from a measured load profile. A plan
written at commissioning describes a village that no longer exists.
\end{principle}

Each group has a different remedy, and the chapters that follow take them in order:
measurement first (Chapter~\ref{ch:measuring-load}), then commissioning and handover
(Chapter~\ref{ch:handover}), which is where knowledge loss is prevented or caused.

\subsection{What this handbook does not cover}

Grid-tied systems, protection design and tariff setting are outside its scope. They
matter, but they are decisions for a qualified engineer, and the cooperative's own
guidance on them is referenced rather than repeated here \cite{coop-protection}.

\section{Measuring load before building}
\label{ch:measuring-load}

Load measurement is the least glamorous part of microgrid work and the part most often
skipped. A two-week logging campaign costs less than one replacement battery module and
answers questions that no survey can: when the evening peak begins, how deep it runs,
and which loads are truly continuous.

\subsection{A two-week logging campaign}

The cooperative logs every feeder at one-minute resolution for fourteen days, spanning
at least one market day and one holiday. Shorter campaigns miss the weekly rhythm;
longer ones rarely change the design decision. Table~\ref{tab:load-summary} summarises
the campaign at the fictional Orrin Ford site.

\begin{table}[h]
\centering
\caption{Fictional fourteen-day load summary for the Orrin Ford site.}
\label{tab:load-summary}
\begin{tabular}{lrrr}
\toprule
Feeder & Mean (kW) & Peak (kW) & Night share \\
\midrule
Clinic & 1.8 & 4.2 & 38\% \\
School & 2.6 & 7.9 & 6\% \\
Market row & 3.1 & 11.4 & 12\% \\
Households & 5.7 & 14.8 & 29\% \\
\midrule
Site total & 13.2 & 31.6 & 22\% \\
\bottomrule
\end{tabular}
\end{table}

\begin{metric}{31.6 kW}{Measured site peak, fourteen-day campaign}
The original design assumed a 22 kW peak. The measured figure is what the battery
and inverter sizing in the next step must serve.
\end{metric}

\subsection{From profile to sizing}

The measured profile feeds three numbers: the inverter's continuous rating, the
battery's usable energy, and the generation needed to refill that energy on a poor
day. The worked method, with the cooperative's safety margins, is in the planning
guide \cite{coop-sizing}; this handbook only insists that the inputs be measured.

\bookpart{Practice}

\section{Commissioning and handover}
\label{ch:handover}

Commissioning is the moment a microgrid is best understood, and handover is the moment
that understanding is most easily lost. The cooperative treats them as one step: a
system is not commissioned until the person who will operate it has run it, alone, for
a week.

\subsection{The handover week}

During the handover week the installer stays on call but does not touch the system.
The incoming operator performs every routine task at least once, records each one in
the site log, and raises every question in writing. Appendix~\ref{app:checklist} lists
the tasks; the glossary defines the terms the log uses.

\begin{itemize}
\item Read and record every meter at the same time each day.
\item Perform one controlled shutdown and restart of the inverter.
\item Walk every feeder and confirm its labels match the single-line diagram.
\item Hand-write the three most likely faults and their first response.
\end{itemize}

\subsection{Keeping the knowledge local}

The site log is the grid's memory. A log that lives on the installer's laptop is not a
log; it is a liability. The cooperative keeps a paper copy on site and a scanned copy
with the district office, and reviews both at every visit \cite{valley-survey}.

\bookappendix

\section{Commissioning checklist}
\label{app:checklist}

The checklist below is the minimum the cooperative accepts before a site is handed
over. Each item is signed by both the installer and the incoming operator.

\begin{enumerate}
\item The measured load profile is attached to the site log.
\item Inverter, battery and generation ratings match the approved sizing sheet.
\item Every feeder is labelled at both ends and on the single-line diagram.
\item A controlled shutdown and restart has been performed by the operator.
\item Earth and insulation test results are recorded and within limits.
\item The three most likely faults and their first responses are posted on site.
\item The district office holds a scanned copy of the complete site log.
\end{enumerate}

\begin{bookglossary}
\glossaryterm{Capacity drift}{The gradual gap between the load a system was sized for
and the load it now serves.}
\glossaryterm{Feeder}{A distribution line leaving the site's main board and serving one
group of loads.}
\glossaryterm{Handover week}{The week in which the incoming operator runs the system
alone while the installer remains on call.}
\glossaryterm{Maintenance debt}{The accumulation of small deferred maintenance tasks,
each harmless on its own.}
\glossaryterm{Site log}{The written record of readings, tasks, faults and questions that
serves as the grid's memory.}
\end{bookglossary}

\begin{bookreferences}{9}
\bibitem{valley-survey} Tessaly Valley Energy Cooperative. \emph{Survey of three
community microgrids, 2023--2026}. Fictional internal report, 2026.
\bibitem{coop-protection} Tessaly Valley Energy Cooperative. \emph{Protection and
earthing guidance for isolated grids}. Fictional guidance note, 2025.
\bibitem{coop-sizing} Tessaly Valley Energy Cooperative. \emph{Planning guide: sizing
storage and generation from measured load}. Fictional planning guide, 2026.
\end{bookreferences}
